\subsection{Reference problem and calculations}

The reference problem follows the published low-Reynolds-number HCCB
pore-scale configuration \cite{wang2023pore,wang2024macro}. Sixty steady
conditions combine five inlet velocities
(\SIrange{0.05}{0.25}{m.s^{-1}}), four inlet temperatures
(\SIrange{300}{900}{K}) and three pebble heat-source levels
(4.85, 6.85 and \SI{8.85}{MW.m^{-3}}). The outlet absolute pressure is
\SI{0.12}{MPa}, the cooled wall is held at \SI{635}{K}, and the spherical
\LiSi{} pebbles have a physical diameter of \SI{1}{mm}. At the inlet,
$Re_{p,\mathrm{in}}=0.078$--$2.40$, $Pr=0.660$--$0.675$ and
$Pe_{p,\mathrm{in}}=0.051$--$1.62$, distinct from the source paper's spatially
averaged $Re_{p,\mathrm{AVE}}<1.8$. The local resolved
domain is cut from the published packing construction and retains the inlet,
outlet, cooled-wall and symmetry-boundary families. A 1\% diameter reduction
is used only to remove point contacts during meshing; consequently, the
reference fields represent an uncompressed bed without solid-contact
conduction.

The steady helium field obeys
\begin{align}
\nabla\cdot(\rho_f\mathbf{u}) &= 0,\\
\rho_f(\mathbf{u}\cdot\nabla)\mathbf{u}
 +\nabla p-\nabla\cdot\!\left[
\mu_f(\nabla\mathbf{u}+\nabla\mathbf{u}^{\mathsf T})\right]&=0,
\end{align}
Here $p$ is absolute thermophysical pressure; \OpenFOAM{} also retains
$p_{rgh}$ for pressure--velocity coupling.
The phase energy equations are
\begin{align}
\rho_fc_{p,f}\left(\frac{\partial T_f}{\partial t}
 +\mathbf{u}\cdot\nabla T_f\right)
 +\nabla\cdot(-k_f\nabla T_f)&=0,\\
\rho_sc_{p,s}\frac{\partial T_s}{\partial t}
 +\nabla\cdot(-k_s\nabla T_s)&=q'''.
\end{align}
Temperature and normal heat flux are continuous at every resolved
fluid--solid interface. Helium density, viscosity and conductivity use the
published HCCB relations, while solid heat storage uses the measured
\LiSi{} enthalpy relation \cite{wang2024macro,kleykamp1996enthalpy}.
The steady calculations set both time derivatives to zero.  In the
fixed-hydrodynamics trajectories, the converged target velocity and pressure
are held constant while the two energy equations remain time dependent.
The published property relations were implemented and verified pointwise in
\OpenFOAM{}, avoiding table extrapolation without implying fully coupled
trajectory accuracy.

All fields are generated with \OpenFOAM{} 13. The steady labels 1--200 are
nonlinear iterations, not physical seconds; accepted cases are finite and
satisfy the reported mass and energy checks.
The 12 paired thermal steps combine the converged source temperature with
target $\mathbf u$, $p$, $p_{rgh}$ and conservative face mass flow $\phi$.
They describe heat transfer after the new flow is established, not the
initial momentum or pressure-wave transient. This follows quasi-steady CHT
methods that exploit fluid--solid time-scale separation
\cite{shen2013quasisteady,gimenez2016weakly,errera2017transientcht}. Its
accuracy during momentum adjustment would require a same-initial-state fully
coupled comparison and is not claimed here. A staged time-step schedule covers
\SIrange{0}{300}{s} and retains 56 common field times per trajectory on the
same fine local mesh.
A representative condition is also solved on coarse, medium and fine fluid--solid meshes
that passed the basic mesh checks. Their unequal equivalent-cell refinement ratios
are retained in the apparent-order and three-grid GCI calculation
\cite{celik2008gci}; percentage GCI is omitted when a near-zero response
crosses sign. Fixed-flow time-step sensitivity is reported separately.

\subsection{Regional graph and physical constraints}

The native finite-volume fields are reduced to \RegionalNodes{} phase-specific
regional nodes connected by \RegionalEdges{} graph edges. The temperature of
region $r$ is the volume average
\begin{equation}
\overline T_r=\frac{\sum_{i\in r}V_iT_i}{\sum_{i\in r}V_i}.
\end{equation}
The volume-weighted projection preserves phase-average temperature and makes
its piecewise-constant error orthogonal to regional prediction error. Thus,
\begin{equation}
E_{T,\mathrm{native}}^2
=E_{T,\mathrm{rep}}^2+E_{T,\mathrm{model}}^2,
\end{equation}
so regional accuracy is not presented as unresolved pore-scale accuracy.

Signed incidence operators assembled from the owner, neighbour and
boundary-face lists of the finite-volume mesh evaluate regional mass and
energy differences:
\begin{equation}
\mathbf b_m=\mathbf A_m\widehat{\dot{\mathbf m}},\qquad
\mathbf b_Q=\mathbf A_Q\widehat{\dot{\mathbf Q}}-\mathbf W.
\end{equation}
The physics-constrained losses compare these quantities with the regionalized
\OpenFOAM{} balances using scales calculated from training conditions only.
Temperature, pressure, interface flux, outlet temperature, pressure drop,
wall heat and hotspot location remain separate reported quantities.

\subsection{Steady and transient models}

The steady comparison uses a response surface, data-only coordinate network,
coordinate PINN, regional graph operator and attention-based physics operator
\cite{raissi2019pinn}. All receive the same conditions and regional fields.
The paired coordinate networks share architecture, initialization and targets;
only the finite-volume mass, energy and interface-flux terms differ.

The transient model composes a frozen steady predictor $P_\psi$, a regional
graph--Transformer $G_\theta$ and an optional diffusion-style temperature
refiner $R_\phi$:
\begin{align}
\widehat{\mathbf z}_0&=
[\widehat{\mathbf u}^{\,\mathrm{tgt}},
\widehat p^{\,\mathrm{tgt}},
\widehat T^{\,\mathrm{src}}],\\
\widehat{\mathbf T}_{1:N_t}
&=G_\theta(\widehat{\mathbf z}_0,
\boldsymbol\xi_{\mathrm{src}},
\boldsymbol\xi_{\mathrm{tgt}},\mathbf t),\\
\mathbf T^{\mathrm{final}}_{1:N_t}
&=\widehat{\mathbf T}_{1:N_t}
+R_\phi(\widehat{\mathbf T}_{1:N_t},\boldsymbol\epsilon).
\end{align}
Training uses the exact \OpenFOAM{} source state to isolate trajectory error.
The strict end-to-end test replaces it with the frozen PINN prediction;
neither network is refitted on the held-out endpoint pair.
Local graph blocks represent short-range exchange, Physics-Attention represents
long-range coupling, and temporal attention advances the trajectory. The
diffusion module changes only temperature and is retained only if it reduces
temperature error without worsening the common energy difference. Persistence,
DMDc and POD corrections provide zero-dynamics, linear-dynamics and low-rank baselines
\cite{vaswani2017attention,proctor2016dmdc,sirovich1987pod,lippe2023pderefiner}.

The common graph--Transformer has 270,441 trainable parameters: a 64-channel
path, two local graph blocks before and after two four-head Physics-Attention
blocks with 128 slices, and three temporal-attention layers. The data-only and
physics-constrained variants share this architecture and initialization; the
factorized variant evaluates the static spatial representation once but keeps
the same temporal queries and parameter count. Training uses at most 500
epochs of AdamW, an initial learning rate of $10^{-3}$, weight decay
$10^{-5}$ and cosine decay to zero. The data-only objective is the
volume-weighted temperature loss $\mathcal L_T$; the constrained variants use
$5\mathcal L_T+\mathcal L_Q+\mathcal L_E$, where $\mathcal L_Q$ compares
regional edge heat flux and $\mathcal L_E$ compares the projected transient
energy operator \cite{li2024pino}. Because fixed coefficients need not balance
the numerical progress of these terms, the predefined $5{:}1{:}1$ baseline is
compared with three published ReLoBRaLo settings \cite{bischof2025relobralo}.
Architecture, split, initialization and training length remain fixed. The
unweighted mean of the three dimensionless validation losses selects both
checkpoint and weighting method; the independent test is read once thereafter.
The loss definition is identical on CPU and CUDA. On one representative
46,089-node regional trajectory, CPU--CUDA and repeated-CUDA residuals differed
by at most $3.0\times10^{-5}$ and $2.7\times10^{-5}$ in relative $L_\infty$
norm, respectively, owing to the order of parallel floating-point summation.

All neural transient models use the same bounded-residual map. For phase $r$,
the network first predicts the unconstrained physical increment
$d_r=(t/t_{\max})s_rq_{\theta,r}$, where $s_r$ is the training-set temperature
scale. Define the available increment in the predicted direction as
\begin{equation}
c_r(d_r)=
\begin{cases}
T_{r,\max}-T_{r,0},&d_r\geq0,\\
T_{r,0}-T_{r,\min},&d_r<0.
\end{cases}
\qquad
\widehat T_r(t)=T_{r,0}+c_r(d_r)
\tanh\!\left[\frac{d_r}{c_r(d_r)}\right].
\label{eq:bounded_temperature_output}
\end{equation}
with \SIrange{300}{1000}{K} for helium and
\SIrange{298}{1300}{K} for the solid. The initial field is exact, and the
limiting value is used when $c_r=0$. The feasible inward derivative is unity at
a bound, so heating from the lower bound and cooling from the upper bound remain
trainable. Equation~\eqref{eq:bounded_temperature_output} enforces the declared
property range without clipping errors within it.

\subsection{Data separation and evaluation}

Steady samples are separated by complete operating condition. The five tests
cover interleaved conditions, high-temperature and high-velocity
extrapolation, and transfer to middle and high heat-source levels. The
source-level tests contain one training source value; its zero-variance input
is therefore zero in every role, making them stability rather than
learned-source-response tests. Transient samples are separated by complete
endpoint pair: both directions joining the
same two steady states remain in one role, and no time point from a held-out
trajectory enters training, normalization or model selection. The strict
end-to-end test also withholds both steady endpoints from steady-model fitting.

Reported metrics include volume-weighted fluid and solid temperature errors,
outlet-temperature and pressure-drop errors, maximum solid temperature,
hotspot displacement, signed wall heat, and regional mass and energy
differences. Transient tests additionally report trajectory error, final-state
error, inference time and diffusion interval calibration. Checkpoints are chosen only
on validation conditions and then frozen for the high-velocity
and independent-packing tests. External PREMUX, TESOMEX, HELOKA and fixed-bed
measurements are used only to check aggregate thermal and hydraulic scales;
they are not cellwise labels for the P418 fields.

Grid and time-step studies quantify numerical sensitivity, the second packing
tests geometric sensitivity, and repeated seeds and training subsets quantify
model and data sensitivity. These effects remain separate and use the same
physical error measures as the final comparison.
